\documentclass[10pt,a4paper]{amsart}
\usepackage{amsmath,amssymb,amsthm,graphicx}
\usepackage{verbatim}
\newenvironment{reference}{\comment}{\endcomment}
%\newenvironment{reference}{}{}
\newenvironment{slogan}{\comment}{\endcomment}
\newenvironment{history}{\comment}{\endcomment}
% For commutative diagrams we use Xy-pic
\usepackage[all]{xy}

% We use 2cell for 2-commutative diagrams.
\xyoption{2cell}
% Theorem environments.
%
\theoremstyle{plain}
\newtheorem{theorem}[subsection]{Theorem}
\newtheorem{proposition}[subsection]{Proposition}
\newtheorem{lemma}[subsection]{Lemma}

\theoremstyle{definition}
\newtheorem{definition}[subsection]{Definition}
\newtheorem{example}[subsection]{Example}
\newtheorem{exercise}[subsection]{Exercise}
\newtheorem{situation}[subsection]{Situation}

\theoremstyle{remark}
\newtheorem{remark}[subsection]{Remark}
\newtheorem{remarks}[subsection]{Remarks}

\numberwithin{equation}{subsection}

% Macros
%
\def\lim{\mathop{\mathrm{lim}}\nolimits}
\def\colim{\mathop{\mathrm{colim}}\nolimits}
\def\Spec{\mathop{\mathrm{Spec}}}
\def\Hom{\mathop{\mathrm{Hom}}\nolimits}
\def\Ext{\mathop{\mathrm{Ext}}\nolimits}
\def\SheafHom{\mathop{\mathcal{H}\!\mathit{om}}\nolimits}
\def\SheafExt{\mathop{\mathcal{E}\!\mathit{xt}}\nolimits}
\def\Sch{\mathit{Sch}}
\def\Mor{\mathop{\mathrm{Mor}}\nolimits}
\def\Ob{\mathop{\mathrm{Ob}}\nolimits}
\def\Sh{\mathop{\mathit{Sh}}\nolimits}
\def\NL{\mathop{N\!L}\nolimits}
\def\CH{\mathop{\mathrm{CH}}\nolimits}
\def\proetale{{pro\text{-}\acute{e}tale}}
\def\etale{{\acute{e}tale}}
\def\QCoh{\mathit{QCoh}}
\def\Ker{\mathop{\mathrm{Ker}}}
\def\Im{\mathop{\mathrm{Im}}}
\def\Coker{\mathop{\mathrm{Coker}}}
\def\Coim{\mathop{\mathrm{Coim}}}

% Boxtimes
%
\DeclareMathSymbol{\boxtimes}{\mathbin}{AMSa}{"02}

%
% Macros for moduli stacks/spaces
%
\def\QCohstack{\mathcal{QC}\!\mathit{oh}}
\def\Cohstack{\mathcal{C}\!\mathit{oh}}
\def\Spacesstack{\mathcal{S}\!\mathit{paces}}
\def\Quotfunctor{\mathrm{Quot}}
\def\Hilbfunctor{\mathrm{Hilb}}
\def\Curvesstack{\mathcal{C}\!\mathit{urves}}
\def\Polarizedstack{\mathcal{P}\!\mathit{olarized}}
\def\Complexesstack{\mathcal{C}\!\mathit{omplexes}}
% \Pic is the operator that assigns to X its picard group, usage \Pic(X)
% \Picardstack_{X/B} denotes the Picard stack of X over B
% \Picardfunctor_{X/B} denotes the Picard functor of X over B
\def\Pic{\mathop{\mathrm{Pic}}\nolimits}
\def\Picardstack{\mathcal{P}\!\mathit{ic}}
\def\Picardfunctor{\mathrm{Pic}}
\def\Deformationcategory{\mathcal{D}\!\mathit{ef}}
\usepackage{cleveref}
\crefname{originalreference}{original reference}{original references}
\Crefname{originalreference}{Original reference}{Original references}
\makeatletter
\newlabel{algebraic-definition-relative-representable-property}{{algebraic-definition-relative-representable-property}{source}{}{sourcepointer-1}{}}
\newlabel{algebraic-definition-relative-representable-property@cref}{{[originalreference][0][]algebraic-definition-relative-representable-property}{[1][0][]source}{}{}{}}
\makeatother
\crefname{originalreference}{original reference}{original references}
\Crefname{originalreference}{Original reference}{Original references}
\makeatletter
\newlabel{algebraic-lemma-characterize-representable-by-algebraic-spaces}{{algebraic-lemma-characterize-representable-by-algebraic-spaces}{source}{}{sourcepointer-2}{}}
\newlabel{algebraic-lemma-characterize-representable-by-algebraic-spaces@cref}{{[originalreference][0][]algebraic-lemma-characterize-representable-by-algebraic-spaces}{[1][0][]source}{}{}{}}
\makeatother
\crefname{originalreference}{original reference}{original references}
\Crefname{originalreference}{Original reference}{Original references}
\makeatletter
\newlabel{criteria-lemma-representable-by-spaces-surjective}{{criteria-lemma-representable-by-spaces-surjective}{source}{}{sourcepointer-3}{}}
\newlabel{criteria-lemma-representable-by-spaces-surjective@cref}{{[originalreference][0][]criteria-lemma-representable-by-spaces-surjective}{[1][0][]source}{}{}{}}
\makeatother
\crefname{originalreference}{original reference}{original references}
\Crefname{originalreference}{Original reference}{Original references}
\makeatletter
\newlabel{criteria-theorem-flat-groupoid-gives-algebraic-stack}{{criteria-theorem-flat-groupoid-gives-algebraic-stack}{source}{}{sourcepointer-4}{}}
\newlabel{criteria-theorem-flat-groupoid-gives-algebraic-stack@cref}{{[originalreference][0][]criteria-theorem-flat-groupoid-gives-algebraic-stack}{[1][0][]source}{}{}{}}
\makeatother
\crefname{originallemma}{original lemma}{original lemmas}
\Crefname{originallemma}{Original lemma}{Original lemmas}
\makeatletter
\newlabel{lemma-quotient-etale}{{lemma-quotient-etale}{source}{}{sourcepointer-5}{}}
\newlabel{lemma-quotient-etale@cref}{{[originallemma][0][]lemma-quotient-etale}{[1][0][]source}{}{}{}}
\makeatother
\crefname{originalreference}{original reference}{original references}
\Crefname{originalreference}{Original reference}{Original references}
\makeatletter
\newlabel{spaces-groupoids-lemma-cartesian-square-of-morphism}{{spaces-groupoids-lemma-cartesian-square-of-morphism}{source}{}{sourcepointer-6}{}}
\newlabel{spaces-groupoids-lemma-cartesian-square-of-morphism@cref}{{[originalreference][0][]spaces-groupoids-lemma-cartesian-square-of-morphism}{[1][0][]source}{}{}{}}
\makeatother
\crefname{originalreference}{original reference}{original references}
\Crefname{originalreference}{Original reference}{Original references}
\makeatletter
\newlabel{spaces-groupoids-lemma-quotient-stack-objects}{{spaces-groupoids-lemma-quotient-stack-objects}{source}{}{sourcepointer-7}{}}
\newlabel{spaces-groupoids-lemma-quotient-stack-objects@cref}{{[originalreference][0][]spaces-groupoids-lemma-quotient-stack-objects}{[1][0][]source}{}{}{}}
\makeatother
\crefname{originalreference}{original reference}{original references}
\Crefname{originalreference}{Original reference}{Original references}
\makeatletter
\newlabel{stacks-properties-lemma-check-property-covering}{{stacks-properties-lemma-check-property-covering}{source}{}{sourcepointer-8}{}}
\newlabel{stacks-properties-lemma-check-property-covering@cref}{{[originalreference][0][]stacks-properties-lemma-check-property-covering}{[1][0][]source}{}{}{}}
\makeatother
\title{A quasi-splitting of a stable vector bundle produces an etale morphism from the quotient}
\author{The Stacks Project Authors}
\date{stacks-project:stacks-morphisms.tex:local-snapshot-20261012; GNU FDL 1.2 or later}
\begin{document}
\maketitle
\noindent\textit{Editorial note.} All delivered mathematical text is the author's own, copied from the pinned Stacks Project chapter source sources/stacks-algebraic-stacks-morphisms/pinned-original-source.tex; no statement or proof step was rewritten, reordered or removed.\par\medskip
\noindent\textit{Editorial note.} Preamble provenance: the Stacks Project keeps one single shared preamble for every chapter; this card takes it directly from the staged sources/stacks-algebraic-stacks-morphisms/preamble.tex as three verbatim spans (lines 8-12, 14-18 and 154-217), so no derived file was written and nothing was added or redefined. The wrapper class and the base packages are supplied by the assembler.\par\medskip
\noindent\textit{Editorial note.} Bibliography provenance: neither the counted statement nor its proof contains a citation (this lane rescans the whole span before writing the card), so the card carries no bibliography and no entry was reconstructed.\par\medskip
\noindent\textit{Editorial note.} Selection provenance: this card was opened under the file-based control contract (CONTROL.md), with the mechanical pre-card gate run on its statement and proof interval before the card was written.\par\medskip
\noindent\textit{Allowed standard prerequisites (editorial).}\begin{itemize}
\item The construction of the quotient stack of the groupoid associated to the quasi-splitting. Used inside the counted proof (source lines 7315-7325). Realises the quotient as the quotient stack of an explicit groupoid.
\item The statement that the relation of the groupoid is an open immersion. Invoked inside the counted proof (source lines 7326-7340). Supplies the hypothesis of the etaleness criterion for flat groupoids.
\item The criterion that a flat finitely presented groupoid with open-immersion relation has an etale quotient. Invoked at the end of the counted proof (source lines 7341-7352). Gives etaleness of the quotient morphism.
\end{itemize}\medskip
\ifdefined\hypertarget\hypertarget{sourcepointer-1}{}\fi
\noindent\textit{Source reference.} Original reference algebraic-definition-relative-representable-property: sources/stacks-algebraic-stacks-morphisms/pinned-original-source.tex, source lines 7332--7332. Navigation pointer only: the counted text refers to the Stacks item algebraic-definition-relative-representable-property, which belongs to another chapter of the project and is not part of this card. The pointer carries the target's own label as its number because this chapter's staged snapshot does not contain the other chapter's tag table; its content is not reproduced here.\par\medskip
\ifdefined\hypertarget\hypertarget{sourcepointer-2}{}\fi
\noindent\textit{Source reference.} Original reference algebraic-lemma-characterize-representable-by-algebraic-spaces: sources/stacks-algebraic-stacks-morphisms/pinned-original-source.tex, source lines 1720--1720. Navigation pointer only: the counted text refers to the Stacks item algebraic-lemma-characterize-representable-by-algebraic-spaces, which belongs to another chapter of the project and is not part of this card. The pointer carries the target's own label as its number because this chapter's staged snapshot does not contain the other chapter's tag table; its content is not reproduced here.\par\medskip
\ifdefined\hypertarget\hypertarget{sourcepointer-3}{}\fi
\noindent\textit{Source reference.} Original reference criteria-lemma-representable-by-spaces-surjective: sources/stacks-algebraic-stacks-morphisms/pinned-original-source.tex, source lines 7335--7335. Navigation pointer only: the counted text refers to the Stacks item criteria-lemma-representable-by-spaces-surjective, which belongs to another chapter of the project and is not part of this card. The pointer carries the target's own label as its number because this chapter's staged snapshot does not contain the other chapter's tag table; its content is not reproduced here.\par\medskip
\ifdefined\hypertarget\hypertarget{sourcepointer-4}{}\fi
\noindent\textit{Source reference.} Original reference criteria-theorem-flat-groupoid-gives-algebraic-stack: sources/stacks-algebraic-stacks-morphisms/pinned-original-source.tex, source lines 3850--3850. Navigation pointer only: the counted text refers to the Stacks item criteria-theorem-flat-groupoid-gives-algebraic-stack, which belongs to another chapter of the project and is not part of this card. The pointer carries the target's own label as its number because this chapter's staged snapshot does not contain the other chapter's tag table; its content is not reproduced here.\par\medskip
\ifdefined\hypertarget\hypertarget{sourcepointer-5}{}\fi
\noindent\textit{Source reference.} Original lemma lemma-quotient-etale: sources/stacks-algebraic-stacks-morphisms/pinned-original-source.tex, source lines 7252--7252. Navigation pointer only: the counted text refers to the Stacks item lemma-quotient-etale, which is an existing labeled result of this same chapter and is not part of this card. The pointer carries the target's own label as its number and its definition line as its span; its content is not reproduced here.\par\medskip
\ifdefined\hypertarget\hypertarget{sourcepointer-6}{}\fi
\noindent\textit{Source reference.} Original reference spaces-groupoids-lemma-cartesian-square-of-morphism: sources/stacks-algebraic-stacks-morphisms/pinned-original-source.tex, source lines 7345--7345. Navigation pointer only: the counted text refers to the Stacks item spaces-groupoids-lemma-cartesian-square-of-morphism, which belongs to another chapter of the project and is not part of this card. The pointer carries the target's own label as its number because this chapter's staged snapshot does not contain the other chapter's tag table; its content is not reproduced here.\par\medskip
\ifdefined\hypertarget\hypertarget{sourcepointer-7}{}\fi
\noindent\textit{Source reference.} Original reference spaces-groupoids-lemma-quotient-stack-objects: sources/stacks-algebraic-stacks-morphisms/pinned-original-source.tex, source lines 7328--7328. Navigation pointer only: the counted text refers to the Stacks item spaces-groupoids-lemma-quotient-stack-objects, which belongs to another chapter of the project and is not part of this card. The pointer carries the target's own label as its number because this chapter's staged snapshot does not contain the other chapter's tag table; its content is not reproduced here.\par\medskip
\ifdefined\hypertarget\hypertarget{sourcepointer-8}{}\fi
\noindent\textit{Source reference.} Original reference stacks-properties-lemma-check-property-covering: sources/stacks-algebraic-stacks-morphisms/pinned-original-source.tex, source lines 479--479. Navigation pointer only: the counted text refers to the Stacks item stacks-properties-lemma-check-property-covering, which belongs to another chapter of the project and is not part of this card. The pointer carries the target's own label as its number because this chapter's staged snapshot does not contain the other chapter's tag table; its content is not reproduced here.\par\medskip
\begin{lemma}
\label{lemma-quasi-splitting-etale}
Let $S$ be a scheme.
Let $(U, R, s, t, c)$ be a groupoid in algebraic spaces over $S$.
Assume $s, t$ are flat and locally of finite presentation.
Let $P \subset R$ be an open subspace such that
$(U, P, s|_P, t|_P, c|_{P \times_{s, U, t} P})$ is a
groupoid in algebraic spaces over $S$. Then
$$
[U/P] \longrightarrow [U/R]
$$
is a morphism of algebraic stacks which
is representable by algebraic spaces, surjective, and \'etale.
\end{lemma}

\begin{proof}
Since $P \subset R$ is open, we see that $s|_P$ and $t|_P$
are flat and locally of finite presentation.
Thus $[U/R]$ and $[U/P]$ are algebraic stacks by
Criteria for Representability, Theorem
\ref{criteria-theorem-flat-groupoid-gives-algebraic-stack}.
To see that the morphism is representable by algebraic spaces,
it suffices to show that $[U/P] \to [U/R]$ is faithful on
fibre categories, see
Algebraic Stacks, Lemma
\ref{algebraic-lemma-characterize-representable-by-algebraic-spaces}.
This follows immediately from the fact that $P \to R$ is a monomorphism
and the explicit description of quotient stacks given in
Groupoids in Spaces, Lemma
\ref{spaces-groupoids-lemma-quotient-stack-objects}.
Having said this, we know what it means for
$[U/P] \to [U/R]$ to be surjective and \'etale by
Algebraic Stacks, Definition
\ref{algebraic-definition-relative-representable-property}.
Surjectivity follows for example from
Criteria for Representability,
Lemma \ref{criteria-lemma-representable-by-spaces-surjective}
and the description of objects of quotient stacks
(see lemma cited above) over spectra of fields.
It remains to prove that our morphism is \'etale.

\medskip\noindent
To do this it suffices to show that $U \times_{[U/R]} [U/P] \to U$
is \'etale, see Properties of Stacks, Lemma
\ref{stacks-properties-lemma-check-property-covering}.
By Groupoids in Spaces, Lemma
\ref{spaces-groupoids-lemma-cartesian-square-of-morphism}
the fibre product is equal to $[R/P \times_{s, U, t} R]$
with morphism to $U$ induced by $s : R \to U$.
The maps $s', t' : P \times_{s, U, t} R \to R$ are given by
$s' : (p, r) \mapsto r$ and $t' : (p, r) \mapsto c(p, r)$. Since
$P \subset R$ is open we conclude that
$(t', s') : P \times_{s, U, t} R \to R \times_{s, U, s} R$
is an open immersion.
Thus we may apply Lemma \ref{lemma-quotient-etale}
to conclude.
\end{proof}

\end{document}
